\subsection{INVERSE LIMITS OF TOPOLOGICAL SPACES}

Let I be a partially ordered (but not necessarily directed) set (*), in which the order relation is written \(\alpha \leqslant \beta\) . For each \(\alpha \in I\) , let \(X_{\alpha}\) be a topological space, and for each pair \((\alpha, \beta)\) such that \(\alpha \leqslant \beta\) let \(f_{\alpha\beta}\) be a mapping of \(X_{\beta}\) into \(X_{\alpha}\) . We say that \((X_{\alpha}, f_{\alpha\beta})\) is an inverse system of topological spaces if: 1) \((X_{\alpha}, f_{\alpha\beta})\) is an inverse system of sets; 2) \(f_{\alpha\beta}\) is a continuous mapping whenever \(\alpha \leqslant \beta\) . Let X denote the set \(\varprojlim X_{\alpha}\) , and for each \(\alpha \in I\) let \(f_{\alpha}\) be the canonical mapping \(X \to X_{\alpha}\) ; then the coarsest topology on X for which the \(f_{\alpha}\) are continuous is said to be the inverse limit (with respect to the \(f_{\alpha\beta}\) ) of the topologies of the \(X_{\alpha}\) , and the set X with this topology is called the inverse limit of the inverse system of topological spaces \((X_{\alpha}, f_{\alpha\beta})\) . Whenever we speak of \(\varprojlim X_{\alpha}\) as a topological space, it is always to be understood that the topology of this space is the inverse limit of the topologies of the \(X_{\alpha}\) , unless the contrary is expressly stated.

The set X is the subset of the product \(\prod_{\alpha\in I}X_{\alpha}\) consisting of those points x such that

\[
\operatorname{pr} _ {\alpha} (x) = f _ {\alpha \beta} (\operatorname{pr} _ {\beta} (x))\tag{I}
\]

whenever \(\alpha \leqslant \beta\) . It follows from Proposition 3 of no. 1 that the inverse limit of the topologies of the \(X_{\alpha}\) is the same as the topology induced on \(X\) by the topology of the product space \(\prod_{\alpha \in I} X_{\alpha}\) . If, for each \(\alpha \in I\) , \(Y_{\alpha}\) is a subspace of \(X_{\alpha}\) such that the \(Y_{\alpha}\) form an inverse system of subsets of the \(X_{\alpha}\) (Set Theory chapter III, § 7, n° 2), then it is clear that the topological space \(\varprojlim Y_{\alpha}\) is a subspace of \(\varprojlim X_{\alpha}\) .

Let \((\mathbf{X}_{\alpha}^{\prime}, f_{\alpha\beta}^{\prime})\) be another inverse system of topological spaces indexed by the same set I, and for each \(\alpha \in I\) let \(u_{\alpha}: X_{\alpha} \to X_{\alpha}^{\prime}\) be a continuous mapping such that \((u_{\alpha})\) is an inverse system of mappings; then \(u = \varprojlim u_{\alpha}\) is a continuous mapping of \(X = \varprojlim X_{\alpha}\) into \(X^{\prime} = \varprojlim X_{\alpha}^{\prime}\) . For if \(f_{\alpha}^{\prime}\) is the canonical mapping \(X^{\prime} \to X_{\alpha}^{\prime}\) , we have \(f_{\alpha}^{\prime} \circ u = u_{\alpha} \circ f_{\alpha}\) , so that \(f_{\alpha}^{\prime} \circ u\) is continuous for each \(\alpha \in I\) , and the assertion follows from Proposition 4 of § 2, no. 3.

Finally, suppose I is a directed set, and let J be a cofinal subset of I; let Z be the inverse limit of the inverse system of topological spaces \((\mathbf{X}_{\alpha}, f_{\alpha\beta})_{\alpha\in\mathbf{J}, \beta\in\mathbf{J}}\) . Then the canonical bijection \(g: X \to Z\) (Set Theory, Chapter III, § 7, no. 2, proposition 3) is a homeomorphism. For we have \(\operatorname{pr}_{\lambda}(g(x)) = \operatorname{pr}_{\lambda}(x)\) for each \(\lambda \in J\) ; hence g is continuous (no. 1, Proposition 1); and if h is the inverse of g, then for each \(\alpha \in I\) there exists \(\lambda \in J\) such that \(\alpha \leqslant \lambda\) , and therefore \(\operatorname{pr}_{\alpha}(h(z)) = f_{\alpha\lambda}(\operatorname{pr}_{\lambda}(z))\) , which shows that h is continuous (no. 1, Proposition 1), since the \(f_{\alpha\lambda}\) are continuous.

PROPOSITION 9. Let I be a directed set and J a cofinal subset of I. Let \((\mathbf{X}_{\alpha}, f_{\alpha \beta})\) be an inverse system of topological spaces indexed by I; let \(\mathbf{X} = \varprojlim \mathbf{X}_{\alpha}\) and let \(f_{\alpha}: \mathbf{X} \to \mathbf{X}_{\alpha}\) be the canonical mapping. Then the family of sets \(\overline{f}_{\alpha}^{1}(\mathbf{U}_{\alpha})\) , where \(\alpha\) runs through J and \(\mathbf{U}_{\alpha}\) runs through a base \(\mathfrak{B}_{\alpha}\) of the topology of \(\mathbf{X}_{\alpha}\) for each \(\alpha \in \mathbf{J}\) , is a base of the topology of \(\mathbf{X}\) .

From § 2, no. 3 we know that the finite intersections of sets of the form \(\overline{f}_{\alpha}^{1}(\mathbf{U}_{\alpha})\) ( \(\alpha \in \mathbf{I}\) , \(\mathbf{U}_{\alpha}\) open in \(\mathbf{X}_{\alpha}\) ) form a base of the topology of \(\mathbf{X}\) . If \((\alpha_i)_{1 \leqslant i \leqslant n}\) is a finite family of indices of \(\mathbf{I}\) , then there exists \(\gamma \in \mathbf{J}\) such that \(\alpha_i \leqslant \gamma\) for \(1 \leqslant i \leqslant n\) ; hence \(f_{\alpha_i} = f_{\alpha_i \gamma} \circ f_\gamma\) ; if we put

\[
\begin{array}{c} \mathrm{V} _ {\gamma} = \bigcap_ {i} \overline {{f}} _ {\alpha_ {i} \gamma} ^ {1} (\mathrm{U} _ {\alpha_ {i}}), \\ \overline {{f}} _ {\gamma} ^ {1} (\mathrm{V} _ {\gamma}) = \bigcap_ {i} \overline {{f}} _ {\alpha_ {i}} ^ {1} (\mathrm{U} _ {\alpha_ {i}}); \end{array}
\]

then we have

but \(V_{\gamma}\) is open and is therefore a union of sets belonging to \(\mathfrak{B}_{\gamma}\) . Hence the result.

COROLLARY. Let A be a subset of X and let \(A_{\alpha}\) denote \(f_{\alpha}(A)\) for each \(\alpha \in I\) . Then:

\begin{enumerate}
\def\labelenumi{(\roman{enumi})}
\item
  The \(\mathbf{A}_{\alpha}\) (resp. the \(\overline{\mathbf{A}}_{\alpha}\) ) form an inverse system of subsets of the \(\mathbf{X}_{\alpha}\) , and \(\overline{\mathbf{A}} = \bigcap_{\alpha} \overline{f}_{\alpha}^{1}(\overline{\mathbf{A}}_{\alpha}) = \varprojlim \overline{\mathbf{A}}_{\alpha}\) .
\item
  If A is closed in X, \(\mathbf{A} = \varprojlim \mathbf{A}_{\alpha} = \varprojlim \overline{\mathbf{A}}_{\alpha}\) .
\end{enumerate}

The first assertion of (i) follows from the relations \(f_{\alpha} = f_{\alpha \beta} \circ f_{\beta}\) for \(\alpha \leqslant \beta\) and from the continuity of the \(f_{\alpha \beta}\) (§ 2, no. 1, Theorem 1). Let \(A'\) denote

\[
\bigcap_ {\alpha} \overline {{f}} _ {\alpha} ^ {1} (\overline {{\mathrm{A}}} _ {\alpha});
\]

then it is clear that \(\mathbf{A}'\) is closed and contains \(\mathbf{A}\) , so that \(\overline{\mathbf{A}} \subset \mathbf{A}'\) . Conversely, let \(x \in \mathbf{A}'\) ; we have to show that \(x\) lies in the closure of \(\mathbf{A}\) . By virtue of Proposition 9 it is enough to prove that every neighbourhood of \(x\) which is of the form \(\overline{f}_{\alpha}^{1}(\mathrm{U}_{\alpha})\) , with \(\alpha \in \mathbf{I}\) and \(\mathbf{U}_{\alpha}\) open in \(\mathbf{X}_{\alpha}\) , meets \(\mathbf{A}\) . Now, by hypothesis, \(f_{\alpha}(x) \in \mathbf{U}_{\alpha}\) , and since \(f_{\alpha}(x) \in \overline{\mathbf{A}}_{\alpha}\) we have \(\mathbf{U}_{\alpha} \cap \mathbf{A}_{\alpha} \neq \emptyset\) , which means that \(\mathbf{A} \cap \overline{f}_{\alpha}^{1}(\mathbf{U}_{\alpha})\) is not empty.

To establish (ii) it is enough to remark that, without any restriction on A, we have \(A \subset \varprojlim A_{\alpha} \subset \varprojlim \overline{A}_{\alpha}\) ; now if A is closed, then from (i)

\[
\mathrm{A} = \varprojlim \overline {{\mathrm{A}}} _ {\alpha}
\]

and (ii) follows.

Example. Let I be a directed set and \((\mathbf{X}_{\alpha})_{\alpha \in I}\) a family of subsets of a set Y, such that \(\mathbf{X}_{\alpha} \supset \mathbf{X}_{\beta}\) whenever \(\alpha \leqslant \beta\) . For each \(\alpha \in I\) let \(\mathfrak{G}_{\alpha}\) be a topology on \(\mathbf{X}_{\alpha}\) such that \(\mathfrak{G}_{\beta}\) is finer than the topology induced on \(\mathbf{X}_{\beta}\) by \(\mathfrak{G}_{\alpha}\) whenever \(\alpha \leqslant \beta\) . If we take \(f_{\alpha \beta}\) to be the canonical injection \(\mathbf{X}_{\beta} \to \mathbf{X}_{\alpha}\) for \(\alpha \leqslant \beta\) , then \(\varprojlim \mathbf{X}_{\alpha}\) may be identified canonically with the intersection X of the \(\mathbf{X}_{\alpha}\) , with the topology which is the least upper bound (§ 2, no. 3, Example 2) of the topologies induced on X by the \(\mathfrak{G}_{\alpha}\) .

